\section{The Benjamini-Hochberg Procedure}
\label{sec:bh}

\subsection{Definition}

\begin{definition}[False Discovery Rate]
Let $R$ be the number of rejected hypotheses and $V$ be the number of
false rejections (false positives).
The false discovery rate is:
\[
  \fdr = E\!\left[\frac{V}{\max(R, 1)}\right].
\]
The FDR is the expected proportion of rejections that are false positives.
\end{definition}

\begin{theorem}[Benjamini-Hochberg Procedure~\citep{bh1995}]
\label{thm:bh}
Let $p_{(1)} \leq p_{(2)} \leq \ldots \leq p_{(m)}$ be the ordered
p-values for $m$ hypothesis tests.
Define the BH threshold:
\[
  k^* = \max\left\{ k : p_{(k)} \leq \frac{k}{m} q \right\}.
\]
Reject all hypotheses $H_{(1)}, H_{(2)}, \ldots, H_{(k^*)}$.
If no such $k^*$ exists, reject nothing.

Under this procedure, $\fdr \leq q$ when the p-values for
the true null hypotheses are independent (or positively correlated).
\end{theorem}

\begin{proof}
See Benjamini and Hochberg~\citep{bh1995}, Theorem 1.
The proof uses a representation of $V/R$ in terms of the spacings of
the null p-values and applies a martingale argument.
$\square$
\end{proof}

\subsection{Properties Relevant to Redistricting}

\paragraph{Power advantage over Bonferroni.}
The Bonferroni correction requires $p_{(k)} \leq \alpha/m$ for any
rejection.
The BH procedure's threshold for the $k$th smallest p-value is
$kq/m$, which is $k$ times larger than Bonferroni for the same $\alpha = q$.
For the ordered hypotheses at the bottom of the list ($k$ close to $m$),
the BH threshold is close to $q$---much less conservative than Bonferroni.

\paragraph{FDR vs.\ FWER.}
Bonferroni controls FWER: $P(\text{any false positive}) \leq \alpha$.
BH controls FDR: $E[\text{fraction of positives that are false}] \leq q$.
FDR control is the appropriate objective when the goal is to identify
a set of significant findings (e.g., ``these states have gerrymandered maps'')
while accepting that a small fraction of the set may be false positives.
FWER control is appropriate when any false positive is catastrophic
(e.g., approving an unsafe drug).

For redistricting, the goal is to identify states with partisan maps
for further legal scrutiny.
A small fraction of false identifications (demoted in subsequent
litigation) is acceptable; missing genuine gerrymanders (false negatives)
is the more costly error.
FDR control at $q = 0.10$ is appropriate.

\paragraph{Positive dependence.}
The BH procedure's FDR guarantee extends to positively correlated
p-values under the Benjamini-Yekutieli (BY) correction~\citep{by2001}.
For redistricting metrics within a single state, the p-values are
positively correlated (a genuinely gerrymandered map tends to be
extreme on multiple metrics simultaneously).
The BY correction would be more conservative than BH; we use BH
as the primary procedure and note that it provides a lower bound
on the FDR guarantee under positive correlation.
